%% =================================================================
%% Chen et al. Measurement 264 (2026) 120300.
%% Reconstruction of page 7 (printed p. 7).
%% Contents: Fig. 9 (a-f: symbol "+" FEA-validation model, FEA
%% simulation of probe pressing, 676-point scan, simulated
%% displacement field, experimental scan result). Prose covers the
%% tail of Section 3.2 (layer-by-layer reconstruction via force
%% thresholds) and Section 4 + 4.1 start.
%% =================================================================

\documentclass[11pt]{article}

\usepackage[margin=0.85in]{geometry}
\usepackage{amsmath, amssymb}
\usepackage{mathrsfs}
\usepackage{xcolor}
\usepackage{fancyhdr}

\pagestyle{fancy}
\fancyhf{}
\lhead{\small Z.~Chen et al.}
\rhead{\small Measurement 264 (2026) 120300}
\cfoot{\small 7 $\to$ \arabic{page}}
\renewcommand{\headrulewidth}{0.4pt}

\setcounter{page}{1}
\setcounter{equation}{1}
\setlength{\parskip}{6pt}
\setlength{\parindent}{0pt}

\begin{document}

% ---------- Fig. 9 placeholder ----------
\begin{center}
\fbox{\begin{minipage}[c][9cm][c]{0.95\textwidth}
\centering
\textbf{Fig. 9.}\ Theoretical validity of the tactile tomography.\\[4pt]
\textcolor{gray}{\footnotesize
\textbf{(a)} Symbol ``\,$+$\,'' model: CAD drawing of a 50 mm $\times$
50 mm substrate with a raised plus-sign feature at the center.\\[3pt]
\textbf{(b)} Symbol ``\,$+$\,'' model filled with AB glues:
photograph showing the glue-covered sample, 50 mm $\times$ 50 mm,
20 mm glue depth marked.\\[3pt]
\textbf{(c)} Simulation results of the mechanical process of the
probe pressing the sample: meshed FEA view with the probe tip (small
yellow sphere) contacting the red sample block.\\[3pt]
\textbf{(d)} Simulation results of the sample scanned by the probe
tip with 676 scanning points: rainbow-colored contour in the
central square region, red-to-yellow elsewhere.\\[3pt]
\textbf{(e)} Simulated displacement field of the sample under probe
loading: rendered 3D surface plot showing the plus-sign structure
emerging via displacement contrast.\\[3pt]
\textbf{(f)} Experimental result of the sample scanned by the
tactile tomography system: cleanly resolved plus-sign with
step-like vertical extrusion on a red substrate (rainbow Z scale).]}
\end{minipage}}
\end{center}

\vspace{0.4em}

\textbf{Fig. 9.}\ Theoretical validity of the tactile tomography.
(a) Symbol ``+'' model. (b) Symbol ``+'' model filled with AB glues.
(c) Simulation results of the mechanical process of the probe
pressing the sample. (d) Simulation results of the sample scanned
by the probe tip with 676 scanning points. (e) Simulated
displacement field of the sample under probe loading. (f)
Experimental result of the sample scanned by the tactile tomography
system.

\vspace{0.8em}

% ---------- Prose ----------
unable to obtain the information of the sample with a thicker soft
surface layer if the applied force is small. As shown in Fig.~8a,
only the top two stages can be reconstructed and distinguished as
the tactile tomography system pressed the sample with a pressure of
0.5 MPa. As the applied force increases, the deeper stage can be
reconstructed and distinguished, such as the top 4 stages for the
pressure of 1.5 MPa (Fig.~8b), the top 6 stages for the pressure of
2.5 MPa (Fig.~8c), and all stage for the pressure of 4 MPa
(Fig.~8d). Therefore, the tactile tomography system can detect and
reconstruct the sample layer-by-layer by setting a series of
threshold corresponding to increasing pressures. For example, nine
thresholds of 0, 1, 2, 3, 4, 5, 6, and 7 that corresponding to the
pressures of 0.5, 0.75, 1.10, 1.50, 1.85, 2.50, 4.00, and 5.00 MPa
were set, and the reconstructed images can be displayed
layer-by-layer in the host computer (as shown in Movie~S1). These
results proved that this imaging system based on the detection of
tactile softness of materials can image an object by tomography.
Moreover, Fig.~8d also shows the maximum depth of 6 mm, indicating
that the tactile tomography system can identify an object buried
6 mm below the surface. The tomography methodology employed in this
tactile tomography system is based on mechanical principles and the
novel parameter of ``tactile softness of materials,'' which
simulates human haptic perception by integrating the effects of
sample size, thickness, and applied force. The system utilizes a
closed-loop feedback control system to drive the probe module to
scan the sample in the X-Y plane with a preset scanning step and
speed. During scanning, the Z-axis rodless cylinder applies a
series of increasing pressures to the sample through the probe. As
the probe interacts with the sample, the pressure sensor detects
force changes, and the grating sensor records the corresponding
compression depth ($\Delta d$) at each X-Y position. This
compression depth, combined with the applied force ($F$) and the
contact area ($S$) between the tip and the sample, defines the
tactile softness (Softness = $\Delta d \cdot S / F$), which
characterizes the material's softness at that location. By
collecting the projection data of positional information
$(x, y, z)$ corresponding to each pressure threshold, the system
reconstructs layer-by-layer 3D images of the internal structure.

\section*{4.\ Results and discussion}

\subsection*{4.1.\ Theoretical validity of the tactile tomography}

To demonstrate the theoretical validity of the tactile tomography,
we performed finite element analysis (FEA) to simulate the
mechanical interaction between the probe tip and samples with
internal features. As shown in Fig.~9a, a model with a feature of a
symbol ``$+$'' was fabricated by 3D printing, and this model was
then filled with a soft surface layer of AB glue (Fig.~9b). A 3D
FEA model was built to simulate the mechanical process of the probe
pressing the sample (Fig.~9c), the sample undergoes

\end{document}
